%% ╔══════════════════════════════════════════════════════════════════╗
%% ║         CORIOTLAB — PLANTILLA HOJA DE VIDA (CV)                 ║
%% ║         Investigadores, docentes, estudiantes de ingeniería     ║
%% ╚══════════════════════════════════════════════════════════════════╝
%%
%% COMPILAR: xelatex -interaction=nonstopmode plantilla_hoja_de_vida.tex
%%           (doble pasada obligatoria)
%%
%% Esta plantilla es autocontenida: no depende de rutas del repositorio
%% "coriotlab-plantillas-informes". Usa los activos de marca ya
%% incluidos en este mismo proyecto (carpetas "705_Kit de Marca
%% Coriotlab" y "Membretes CoriotLab", un nivel arriba de este archivo).

%% ═══════════════════════════════════════════════════════════════════
%% ZONA DE CONFIGURACIÓN — únicos campos a editar
%% ═══════════════════════════════════════════════════════════════════

%% --- Identificación ---
\newcommand{\CVNombre}{Nombre Apellido}
\newcommand{\CVTitulo}{Ingeniero(a) Electrónico(a) / Investigador(a) en Robótica}
\newcommand{\CVCorreo}{correo@ejemplo.com}
\newcommand{\CVTelefono}{+57 300 000 0000}
\newcommand{\CVCiudad}{Medellín, Colombia}
\newcommand{\CVLinkedIn}{}       %% opcional — URL completa, ej. https://www.linkedin.com/in/usuario
\newcommand{\CVORCID}{}         %% opcional — solo el identificador, ej. 0000-0000-0000-0000
\newcommand{\CVGitHub}{}        %% opcional — URL completa, ej. https://github.com/usuario
\newcommand{\CVGoogleScholar}{} %% opcional — URL completa del perfil de Google Scholar
\newcommand{\CVResearchGate}{}  %% opcional — URL completa del perfil de ResearchGate
%% Deje cualquiera de los anteriores vacío ({}) si no aplica — el CV los omite automáticamente.
\newcommand{\CVFecha}{Julio de 2026}   %% fecha de última actualización del CV

%% ═══════════════════════════════════════════════════════════════════

\documentclass[10pt, letterpaper]{article}

%% --- Idioma
\usepackage[spanish, es-nodecimaldot, es-noshorthands]{babel}

%% --- Fuentes (nombres de familia Windows — instalar con setup_coriotlab.ps1)
\usepackage{fontspec}
\setmainfont{Inter}
\newfontfamily\FuenteTitulo{MuseoModerno}
\newfontfamily\FuenteCodigo{Space Mono}

%% --- Geometría — más compacta que los informes: el CV no usa membrete
%%     de página completa, así que se aprovecha más ancho de contenido.
\usepackage[
  top=2.1cm, bottom=2.3cm,
  left=2.2cm, right=2.2cm,
  footskip=1.1cm
]{geometry}

%% --- Colores corporativos CORIOTLAB (idénticos al Libro de Marca)
\usepackage[table]{xcolor}
\definecolor{AzulITM}     {HTML}{102D69}
\definecolor{Magenta}     {HTML}{C14894}
\definecolor{AzulDigital} {HTML}{56ACDE}
\definecolor{GrisPizarra} {HTML}{2F2F2F}
\definecolor{GrisClaro}   {HTML}{F5F5F5}
\definecolor{GrisMedio}   {HTML}{888888}
\definecolor{GrisLinea}   {HTML}{DDDDDD}
\definecolor{Blanco}      {HTML}{FFFFFF}
\definecolor{VerdeTarea}  {HTML}{1A6B3C}

%% --- Gráficos — grffile permite rutas con espacios y puntos múltiples,
%%     necesario porque las carpetas del kit de marca usan ese formato
%%     (ej. "2. Logos/2. Logo Full Color/2. Horizontal/...")
\usepackage{graphicx}
\usepackage{grffile}
\usepackage{tikz}

%% --- Tablas y alineación de entradas
\usepackage{tabularx}
\usepackage{array}
\usepackage{colortbl}
\usepackage{booktabs}

%% --- Listas
\usepackage{enumitem}

%% --- Campos opcionales condicionales (DOI, LinkedIn, ORCID, GitHub, etc.)
\usepackage{etoolbox}

%% --- Secciones y espaciado
\usepackage{titlesec}
\usepackage{needspace}
\usepackage[hidelinks]{hyperref}
\hypersetup{urlcolor=AzulDigital}

%% --- Pie de página discreto (fecha + numeración si el CV excede 1 página)
\usepackage{fancyhdr}

\usepackage{parskip}
\setlength{\parskip}{4pt}
\setlength{\parindent}{0pt}
\renewcommand{\arraystretch}{1.3}

%% ─────────────────────────────────────────────────────────────────
%% PIE DE PÁGINA — regla delgada + fecha (izq) + página (der)
%% Reemplaza el membrete de página completa que usan los informes:
%% en una hoja de vida (densa en tablas/listas) un watermark central
%% compite con el contenido y reduce el contraste en blanco y negro.
%% ─────────────────────────────────────────────────────────────────
\pagestyle{fancy}
\fancyhf{}
\renewcommand{\headrulewidth}{0pt}
\renewcommand{\footrulewidth}{0.4pt}
\renewcommand{\footrule}{{\color{GrisLinea}\hrule height\footrulewidth}}
\fancyfoot[L]{{\FuenteCodigo\tiny\color{GrisMedio}\CVFecha}}
\fancyfoot[R]{{\small\color{GrisMedio}\thepage}}

%% ─────────────────────────────────────────────────────────────────
%% ESTILO DE SECCIÓN — cuadrado de acento (píxeles del logo) + título
%% MuseoModerno en AzulITM + regla AzulDigital, sin numeración (no es
%% un informe con TOC, es un documento de una sola lectura corrida).
%% ─────────────────────────────────────────────────────────────────
\newcommand{\SeccionCV}[1]{%
  \needspace{3\baselineskip}%
  \vspace{1.0em}%
  \noindent{\textcolor{AzulITM}{\rule{2.4mm}{2.4mm}}}\hspace{0.5em}%
  {\FuenteTitulo\large\color{AzulITM}\bfseries #1}\\[0.2em]%
  {\color{AzulDigital}\rule{\linewidth}{1.1pt}}%
  \vspace{0.4em}%
}

%% ─────────────────────────────────────────────────────────────────
%% ENCABEZADO DE ENTRADA — fila título (izq) + fechas (der), usado
%% en Formación, Experiencia, Certificaciones. Mantiene alineación
%% consistente sin obligar a usar una tabla completa por sección.
%% ─────────────────────────────────────────────────────────────────
\newcommand{\ItemCV}[2]{%
  \par\noindent\begin{tabularx}{\linewidth}{@{}X r@{}}
    {\FuenteTitulo\normalsize\color{GrisPizarra}\bfseries #1} &
    {\FuenteCodigo\small\color{GrisMedio}#2}\\
  \end{tabularx}\par\vspace{0.1em}%
}

%% Estilo común para listas con viñeta "<" del sistema de marca (3.4).
%% align=left evita que enumitem intente encajar la etiqueta en un
%% labelwidth insuficiente y la empuje fuera del margen izquierdo.
\setlist[itemize]{label={\color{AzulITM}\textbf{>}},
  leftmargin=1.6em, labelsep=0.5em, align=left, itemsep=2pt, topsep=3pt}

%% ─────────────────────────────────────────────────────────────────
%% SISTEMA DE ETIQUETAS — Nivel de formación académica (sección 4A)
%% ─────────────────────────────────────────────────────────────────
\newcommand{\NivelFormacion}[1]{%
  \colorbox{AzulITM!12}{\color{AzulITM}\small\bfseries\strut\ #1\ }}
\newcommand{\FormacionEnCurso}{%
  \colorbox{AzulDigital!15}{\color{AzulDigital!70!black}\small\bfseries\strut\ En curso\ }}
\newcommand{\FormacionCulminada}{%
  \colorbox{VerdeTarea!15}{\color{VerdeTarea}\small\bfseries\strut\ Culminado\ }}

%% ─────────────────────────────────────────────────────────────────
%% SISTEMA DE ETIQUETAS — Certificación (sección 4B)
%% ─────────────────────────────────────────────────────────────────
\newcommand{\Certificado}{%
  \colorbox{VerdeTarea!15}{\color{VerdeTarea}\small\bfseries\strut\ Certificado\ }}
\newcommand{\NoCertificado}{%
  \colorbox{GrisMedio!15}{\color{GrisMedio!70!black}\small\bfseries\strut\ Sin certificar\ }}

%% ─────────────────────────────────────────────────────────────────
%% PUBLICACIONES EN FORMATO APA 7 (Opción A — ver README, sección
%% "Decisión técnica: publicaciones en APA"). Comando determinístico,
%% sin biblatex/biber ni archivo .bib externo.
%%
%% Uso:
%%   \Publicacion{Autores}{Año}{Título}{Fuente en cursiva}%
%%               {Volumen(Número), páginas}{DOI o URL}
%%
%% Deje {} vacío en el 5.º o 6.º parámetro si no aplica (ej. libros
%% sin volumen, o publicaciones sin DOI).
%% El orden de aparición es manual: para reordenar (cronológico o
%% alfabético) simplemente reordene las llamadas a \Publicacion en
%% el documento.
%% ─────────────────────────────────────────────────────────────────
\newcommand{\Publicacion}[6]{%
  \par\noindent\hangindent=1.2em\hangafter=1
  #1 (#2). #3. \textit{#4}%
  \ifstrempty{#5}{}{, \textit{#5}}%
  .%
  \ifstrempty{#6}{}{\ {\FuenteCodigo\small #6}}%
  \par\vspace{0.5em}%
}

%% ─────────────────────────────────────────────────────────────────
%% CONGRESOS Y PONENCIAS — formato de cita breve, visualmente
%% alineado con \Publicacion pero sin exigir rigor APA estricto.
%%
%% Uso:
%%   \Ponencia{Título de la ponencia}{Nombre del evento}%
%%            {Tipo de participación}{Lugar}{Fecha}
%% ─────────────────────────────────────────────────────────────────
\newcommand{\Ponencia}[5]{%
  \par\noindent\hangindent=1.2em\hangafter=1
  \textbf{#1}. \textit{#2}. #3, #4 (#5).
  \par\vspace{0.5em}%
}

%% =================================================================
%%  DOCUMENTO
%% =================================================================
\begin{document}
\color{GrisPizarra}

%% ─────────────────────────────────────────────────────────────────
%% ENCABEZADO — nombre (protagonista) + logo horizontal discreto
%% El logo va en menor jerarquía visual que el nombre del titular,
%% tal como exige el Libro de Marca (sección "Usos correctos").
%% ─────────────────────────────────────────────────────────────────
\noindent
\begin{minipage}[t]{0.72\linewidth}
  {\FuenteTitulo\fontsize{24}{28}\selectfont\color{AzulITM}\bfseries \CVNombre}\\[0.15em]
  {\large\color{GrisPizarra}\CVTitulo}
\end{minipage}%
\hfill
\begin{minipage}[t]{0.26\linewidth}
  \raggedleft
  \includegraphics[width=0.85\linewidth]%
    {../705_Kit de Marca Coriotlab/2. Logos/2. Logo Full Color/2. Horizontal/Coriotlab_Logo_Hor_FullColor.png}
\end{minipage}

\vspace{0.6em}

%% --- Línea de contacto: correo, teléfono, ciudad ---
\newcommand{\CVSep}{\ \ \textcolor{AzulITM}{\textbf{<}}\ \ }
\noindent{\small\color{GrisMedio}
  \CVCorreo
  \ifdefempty{\CVTelefono}{}{\CVSep\CVTelefono}%
  \CVSep\CVCiudad
}

%% --- Línea de redes/identificadores académicos (separada del contacto) ---
%% Solo aparece si al menos uno de estos campos no está vacío.
\newbool{CVMostrarRedes}
\ifdefempty{\CVLinkedIn}{}{\booltrue{CVMostrarRedes}}
\ifdefempty{\CVORCID}{}{\booltrue{CVMostrarRedes}}
\ifdefempty{\CVGitHub}{}{\booltrue{CVMostrarRedes}}
\ifdefempty{\CVGoogleScholar}{}{\booltrue{CVMostrarRedes}}
\ifdefempty{\CVResearchGate}{}{\booltrue{CVMostrarRedes}}
\ifbool{CVMostrarRedes}{%
  \vspace{0.25em}

  \noindent{\small\color{GrisMedio}\textit{Redes:}\ \,%
    \def\CVPrimeraRed{1}%
    \ifdefempty{\CVLinkedIn}{}{\href{\CVLinkedIn}{LinkedIn}\def\CVPrimeraRed{0}}%
    \ifdefempty{\CVORCID}{}{\ifnum\CVPrimeraRed=0\CVSep\fi\href{https://orcid.org/\CVORCID}{{\FuenteCodigo ORCID: \CVORCID}}\def\CVPrimeraRed{0}}%
    \ifdefempty{\CVGoogleScholar}{}{\ifnum\CVPrimeraRed=0\CVSep\fi\href{\CVGoogleScholar}{Google Scholar}\def\CVPrimeraRed{0}}%
    \ifdefempty{\CVResearchGate}{}{\ifnum\CVPrimeraRed=0\CVSep\fi\href{\CVResearchGate}{ResearchGate}\def\CVPrimeraRed{0}}%
    \ifdefempty{\CVGitHub}{}{\ifnum\CVPrimeraRed=0\CVSep\fi\href{\CVGitHub}{GitHub}\def\CVPrimeraRed{0}}%
  }%
}{}

\vspace{0.3em}
{\color{GrisLinea}\rule{\linewidth}{0.6pt}}

%% ── OPT-0: Perfil Profesional ─────────────────────────────────────
%% Opcional (sección 2, ítem 3 del encargo). Descomente todo el bloque
%% para activarlo.
%\SeccionCV{Perfil Profesional}
%[Párrafo breve de 3 a 5 líneas: área de especialización, enfoque de
%investigación o línea de trabajo principal, y el tipo de proyectos o
%resultados que busca destacar. Evite repetir aquí lo que ya se detalla
%en Experiencia o Formación; use este espacio para dar contexto y
%orientación al lector sobre su perfil general.]

%% ─────────────────────────────────────────────────────────────────
%% FORMACIÓN ACADÉMICA
%% ─────────────────────────────────────────────────────────────────
\SeccionCV{Formación Académica}

\ItemCV{[Nombre del programa académico]}{[AAAA] -- en curso}
[Institución]\quad\NivelFormacion{[Pregrado / Especialización / Maestría / Doctorado]}\quad\FormacionEnCurso

\vspace{0.7em}

\ItemCV{[Nombre del programa académico]}{[AAAA] -- [AAAA]}
[Institución]\quad\NivelFormacion{[Nivel]}\quad\FormacionCulminada

%% ── OPT-A: Idiomas ───────────────────────────────────────────────
%\SeccionCV{Idiomas}
%\begin{itemize}
%  \item \textbf{Español} --- Nativo
%  \item \textbf{Inglés} --- [Nivel: A1--C2 o descriptor MCER]
%\end{itemize}

%% ─────────────────────────────────────────────────────────────────
%% EXPERIENCIA CERTIFICADA
%% ─────────────────────────────────────────────────────────────────
\SeccionCV{Experiencia Certificada}

\ItemCV{[Cargo o rol desempeñado]}{[Mes AAAA] -- [Mes AAAA]}
[Empresa o institución]\quad\Certificado
\begin{itemize}
  \item{} [Función o logro concreto, con verbo en pasado y resultado
        verificable cuando sea posible.]
  \item{} [Segunda función o logro relevante.]
\end{itemize}

\vspace{0.5em}

\ItemCV{[Cargo o rol desempeñado]}{[Mes AAAA] -- Actualidad}
[Empresa o institución]\quad\Certificado
\begin{itemize}
  \item{} [Función o logro concreto.]
  \item{} [Segunda función o logro relevante.]
\end{itemize}

%% ─────────────────────────────────────────────────────────────────
%% CERTIFICACIONES
%% ─────────────────────────────────────────────────────────────────
\SeccionCV{Certificaciones}

\ItemCV{[Nombre de la certificación o curso]}{[Mes AAAA]}
[Entidad emisora]\quad{\FuenteCodigo\small\color{GrisMedio}[Código de verificación o enlace, si aplica]}

\vspace{0.5em}

\ItemCV{[Nombre de la certificación o curso]}{[Mes AAAA]}
[Entidad emisora]\quad{\FuenteCodigo\small\color{GrisMedio}[Código de verificación o enlace, si aplica]}

%% ─────────────────────────────────────────────────────────────────
%% PUBLICACIONES (formato APA 7 — ver comando \Publicacion arriba)
%% ─────────────────────────────────────────────────────────────────
\SeccionCV{Publicaciones}

\Publicacion{Apellido, N. A., y Apellido, N. B.}{AAAA}%
  {Título del artículo en formato de oración normal}%
  {Nombre de la Revista}{Vol(Núm), pp--pp}%
  {https://doi.org/xx.xxxx/xxxxxxx}

\Publicacion{Apellido, N. A.}{AAAA}%
  {Título de la segunda publicación}%
  {Nombre de la Revista o Editorial}{}{}

%% ─────────────────────────────────────────────────────────────────
%% CONGRESOS Y PONENCIAS
%% ─────────────────────────────────────────────────────────────────
\SeccionCV{Congresos y Ponencias}

\Ponencia{Título de la ponencia o trabajo presentado}%
  {Nombre del congreso o evento}%
  {Ponente}{Ciudad, País}{AAAA}

\Ponencia{Título del trabajo presentado}%
  {Nombre del congreso o evento}%
  {Póster}{Ciudad, País}{AAAA}

%% ─────────────────────────────────────────────────────────────────
%% SECCIONES OPCIONALES ADICIONALES — descomente la que aplique
%% (sección 2, ítem 9 del encargo: no fueron pedidas explícitamente,
%% pero son estándar en un CV técnico-investigativo)
%% Nota: OPT-A (Idiomas) está ubicado más arriba, entre Formación
%% Académica y Experiencia Certificada.
%% ─────────────────────────────────────────────────────────────────

%% ── OPT-B: Reconocimientos ───────────────────────────────────────
%% Premios, distinciones o participación como evaluador/jurado.
%\SeccionCV{Reconocimientos}
%\ItemCV{Nombre del reconocimiento o premio}{Mes AAAA}
%[Entidad otorgante]
%\begin{itemize}
%  \item{} [Proyecto o motivo del reconocimiento, coautores si aplica.]
%\end{itemize}

%% ── OPT-C: Habilidades técnicas / Stack tecnológico ──────────────
%\SeccionCV{Habilidades Técnicas}
%{\small
%\textbf{Lenguajes:} [Python, C++, MATLAB, \ldots] \\
%\textbf{Frameworks / Librerías:} [ROS, TensorFlow, \ldots] \\
%\textbf{Hardware / Plataformas:} [Arduino, Raspberry Pi, PLC, \ldots] \\
%\textbf{Herramientas de laboratorio:} [Osciloscopio, analizador lógico, \ldots]
%}

%% ── OPT-D: Software y herramientas ───────────────────────────────
%\SeccionCV{Software y Herramientas}
%\begin{itemize}
%  \item{} [Nombre del software] --- [Nivel de dominio o uso específico]
%\end{itemize}

\end{document}
